\section{Learners and Baselines}\label{sec:learners}

All step policies are Gaussian MLPs (64--64 units for tracking, 128--64 for navigation) with a state-independent standard deviation. To keep differences down to the meta-learning rule, the gradient-based methods share one estimator: per-task linear feature baselines with generalised advantage estimation~\cite{schulman2015gae} and PPO-clipped updates~\cite{schulman2017ppo}.

\textbf{MAML family.} Each iteration samples 16 tasks and flies 10 episodes per task. One vanilla policy-gradient step with step size 0.1 adapts the policy to each task; all tasks adapt in one batched backward pass. The adapted policies fly 10 new episodes, and a PPO-clipped outer loss on those episodes is backpropagated through the adaptation step. MAML~\cite{finn2017maml} differentiates through the inner gradient (second order); FOMAML drops that term; ANIL adapts only the output layer and the standard deviation~\cite{raghu2020anil}; Meta-SGD also learns a step size for every parameter~\cite{li2017metasgd}. Like the original MAML for RL, we ignore how the pre-adaptation sampling distribution depends on the parameters~\cite{rothfuss2018promp}.

\textbf{Reptile}~\cite{nichol2018reptile} runs three rounds of per-task PPO (Adam, four epochs per round), each on fresh episodes, then moves the shared parameters toward the mean adapted parameters with a step that anneals from 0.5 to 0.1.

\textbf{PEARL}~\cite{rakelly2019pearl} trains a soft actor-critic~\cite{haarnoja2018sac} conditioned on a five-dimensional latent task variable. A permutation-invariant encoder maps transitions from 32 training tasks to a product-of-Gaussians posterior, trained through the critic loss plus a KL penalty to the prior.

\textbf{RL$^2$}~\cite{duan2016rl2,wang2016learning} uses a GRU policy with 128 units that sees the previous action, reward and an episode-start flag, and keeps its hidden state over a trial of three episodes on the same task. It is trained with recurrent PPO over whole trials.

\textbf{Baselines.} \emph{Classical} uses nominal gains and no learning; \emph{Classical + L1} adds L1 augmentation. \emph{DR} trains one policy with PPO over the whole task distribution (domain randomisation) and never adapts; \emph{DR + fine-tune} adapts it at test time with the same per-task PPO rounds as Reptile. For Combination 2 the DR baseline is the same Gaussian over gains, trained over all tasks without a meta objective. For navigation we add \emph{end-to-end DR}, a DR planner that sees the goal direction instead of the carrot, with no map or planner.

\textbf{Training budgets.} Table~\ref{tab:budget} lists the training budget of each learner. DR and the MAML family saw the same data on tracking, while PEARL and RL$^2$ saw less.

\begin{table}[t]
\centering
\caption{Training budget in millions of environment steps.}
\label{tab:budget}
\footnotesize
% generated by figures_scripts/budget_table.py; do not edit
\begin{tabular}{@{}lrrrrr@{}}
\toprule
Combination & DR & MAML fam. & Reptile & PEARL & RL$^2$ \\
\midrule
C1 residual & 9.6 & 9.6 & 11.5 & 1.0 & 4.3 \\
C2 gains & 4.8 & 6.4 & 7.7 & -- & -- \\
C3 planner$^*$ & 12.0 & 5.8 & 7.2 & 1.0 & 5.2 \\
\bottomrule
\multicolumn{6}{@{}l}{$^*$MAML family and Reptile: additional steps after starting from DR.}
\end{tabular}

\end{table}

\textbf{Navigation specifics.} Training navigation planners from scratch is the most expensive part of the study, so to fit our CPU budget the gradient-based meta-learners start from the trained DR planner and meta-train from there. PEARL and RL$^2$ use different networks and train from scratch, so their comparison with the others is confounded by starting point. With the tracking step size (Adam, 3e-3), per-task PPO adaptation wrecked the warm-started navigation planner (Reptile reached 0\% success), so navigation uses 3e-4 for both Reptile and DR + fine-tune; the diverged run is kept in the released results.
